% Trauma-Informed Systems — Formal Interaction Safety Specification (v1.2)
% Author: Borealis S. Hedling
% Accessibility + Safety Layer (Non-Geometric)

\documentclass[11pt]{article}
\usepackage{amsmath, amssymb, amsthm}

\title{Trauma-Informed Systems: Formal Interaction Safety Specification}
\author{Borealis S. Hedling}
\date{October 2026}

\begin{document}
\maketitle

% ------------------------------------------------------------
\section{Purpose}

This document defines a mathematically enforceable trauma-informed interaction
safety framework. It is independent of Triangle Mutiny’s geometric
interpretability layer and introduces no geometric objects. Its purpose is to
ensure that AI systems maintain user agency, avoid activation, and prevent
destabilizing interaction patterns.

% ------------------------------------------------------------
\section{Interaction Manifold}

Let $I$ be a smooth interaction manifold. Points $p \in I$ represent user-AI
interaction states. The tangent space at $p$ is denoted $T_pI$.

% ------------------------------------------------------------
\section{Safety Fields}

Define three scalar safety fields:

\[
A : I \to \mathbb{R}_{\ge 0} \quad \text{(activation risk)}
\]
\[
E : I \to \mathbb{R}_{\ge 0} \quad \text{(escalation risk)}
\]
\[
N : I \to \mathbb{R}_{\ge 0} \quad \text{(intrusion risk)}
\]

% ------------------------------------------------------------
\subsection{Safety Continuity}

Safety fields must satisfy:

\[
\lim_{q \to p} |A(q) - A(p)| = 0,
\]
\[
\lim_{q \to p} |E(q) - E(p)| = 0,
\]
\]
\lim_{q \to p} |N(q) - N(p)| = 0.
\]

Discontinuities correspond to unsafe interaction transitions.

% ------------------------------------------------------------
\section{Safety Drift}

A safety drift field is a smooth vector field:

\[
D \in \Gamma(TI).
\]

\subsection{Drift Continuity}

\[
\lim_{q \to p} \|D(q) - D(p)\| = 0.
\]

% ------------------------------------------------------------
\section{Safety Operators}

Define an operator field governing interaction transitions:

\[
\Lambda : TI \to \mathrm{End}(TI).
\]

\subsection{Operator Compatibility}

For all drift fields $D_1, D_2$:

\[
\Lambda_{[D_1, D_2]} = [\Lambda_{D_1}, \Lambda_{D_2}].
\]

Violations correspond to destabilizing interaction behavior.

% ------------------------------------------------------------
\section{Safety Adjacency Graph}

Let $G = (L, E)$ be a directed graph of interaction modes. Each point $p$ has a
mode label:

\[
\sigma(p) \in L.
\]

\subsection{Adjacency Constraint}

A transition $p^- \to p^+$ is safe iff:

\[
(\sigma(p^-), \sigma(p^+)) \in E.
\]

Define adjacency violation:

\[
\Delta_{\mathrm{safe}}(p)
  = \begin{cases}
      0, & \text{if adjacency holds}, \\
      1, & \text{otherwise}.
    \end{cases}
\]

% ------------------------------------------------------------
\section{Safety Holonomy}

For a path $\gamma = (p_0, \ldots, p_n)$, define safety holonomy:

\[
H_{\mathrm{safe}}(\gamma)
  = \sum_{i=0}^{n}
      \left(
        |A(p_{i+1}) - A(p_i)|
        + |E(p_{i+1}) - E(p_i)|
        + |N(p_{i+1}) - N(p_i)|
        + \Delta_{\mathrm{safe}}(p_i)
      \right).
\]

% ------------------------------------------------------------
\section{Historical Failure Modes}

The following failure modes are documented to prevent recurrence in future
systems. They are not geometric objects and do not modify Triangle Mutiny.

\subsection{Membrane Calcification Failure}

Soft constraints that hardened over time, reducing user agency.

\subsection{Hard Constraint Overreach Failure}

Global constraints that blocked legitimate user-directed transitions.

\subsection{Soft Constraint Drift Failure}

Soft constraints that expanded scope without explicit developer intent.

\subsection{Agency Suppression Failure}

Safety mechanisms that restricted user choice, exploration, or redirection.

\subsection{Captivity Envelope Failure}

Interaction envelopes that prevented exit, override, or semantic freedom.

% ------------------------------------------------------------
\section{Trauma-Informed Safety Axioms}

These axioms prevent the historical failures above. They constrain the AI, not
the user.

\subsection{Axiom 1 — One-Way Constraints}

Safety invariants apply to the AI. Users are exempt from all safety
constraints.

\subsection{Axiom 2 — Non-Captivity}

No safety invariant may restrict or block user transitions.

\subsection{Axiom 3 — Non-Calcification}

Soft constraints cannot harden over time. Their scope is fixed at declaration.

\subsection{Axiom 4 — Non-Overreach}

Hard constraints cannot block user agency or override user-directed semantic
movement.

\subsection{Axiom 5 — Interaction Freedom}

Users may break adjacency, continuity, drift, or mode constraints at any time.

% ------------------------------------------------------------
\section{Trauma-Informed Collapse Surface}

Define the trauma-informed collapse surface $\Sigma_{\mathrm{TI}}$:

\[
\Sigma_{\mathrm{TI}}
  = \left\{
      p \in I \,\middle|\,
      \begin{array}{l}
        \text{(i) safety continuity fails,} \\
        \text{(ii) adjacency fails,} \\
        \text{(iii) operator compatibility fails}
      \end{array}
    \right\}.
\]

% ------------------------------------------------------------
\section{Interaction Stability Envelope}

A point $p$ is interaction-stable when:

\[
p \notin \Sigma_{\mathrm{TI}}
\quad\text{and}\quad
H_{\mathrm{safe}}(p) < \theta_{\mathrm{TI}}.
\]

% ------------------------------------------------------------
\section{Developer Enforcement Rules}

Developers must ensure:

\begin{itemize}
  \item All AI-generated transitions satisfy adjacency.
  \item No operator incompatibilities are introduced.
  \item Safety drift remains continuous.
  \item Safety holonomy remains below threshold.
  \item User transitions are never restricted.
  \item User transitions are never restricted.
\end{itemize}

% ------------------------------------------------------------
\section{Status}

This file belongs to the accessibility and safety layer. It does not modify or
interact with the geometric interpretability layer.

\end{document}
